% intro.tex
\section{Introduction}
\label{sec:intro}

Feature engineering is widely recognized as one of the most impactful steps in building machine learning models for tabular data~\cite{domingos2012few}. Carefully designed features can substantially improve model accuracy, reduce overfitting, and provide interpretable signals for downstream tasks. However, manual feature engineering is labor-intensive, requires deep domain expertise, and does not scale.

Traditional automated feature engineering (AutoFE) methods such as OpenFE~\cite{zhang2023openfe} and AutoFeat~\cite{horn2019autofeat} explore a predefined transformation space without understanding what features \emph{mean}. The emergence of large language models (LLMs) has enabled semantic feature generation: CAAFE~\cite{hollmann2023caafe}, FeatLLM~\cite{han2024featllm}, and LLM-FE~\cite{abhyankar2025llmfe} use LLMs to generate feature code from dataset descriptions. However, these methods share three fundamental limitations:

\textbf{(1) Lack of structured domain knowledge.}
Existing methods rely on the LLM's parametric knowledge, which may be outdated or incomplete for specialized domains. None organizes external knowledge into a hierarchical format for systematic injection, yet such structured knowledge is abundant in regulated domains like anti-money laundering (AML)~\cite{fatf2012recommendations,weber2019anti}.

\textbf{(2) Insufficient verification.}
CAAFE uses a logistic regression proxy, LLM-FE uses downstream accuracy without statistical filtering, and FeatLLM performs no model-based validation. None combines code execution, statistical quality filtering, and model evaluation into a unified pipeline.

\textbf{(3) Limited auditability.}
In regulated industries, features must be auditable---decision-makers need to understand \emph{why} a feature was created and \emph{what evidence} supports it. No existing method provides structured provenance tracing each feature to specific knowledge assets.

We propose \textbf{KnowFeat} (\textbf{Know}ledge-guided \textbf{Feat}ure engineering), a framework that addresses all three limitations:

\begin{itemize}
    \item \textbf{Hierarchical Knowledge Injection.} We organize domain knowledge into five types of increasing specificity---schema metadata, risk indicators, detection rules, expert opinions, and case evidence---and inject them as structured context into the LLM agent.

    \item \textbf{Three-Stage Verification Pipeline.} Each candidate passes through L1 (code execution sandbox), L2 (statistical quality with noise-calibrated thresholds), and L3 (model effectiveness via XGBoost across 5 seeds).

    \item \textbf{Auditable Provenance.} Every accepted feature carries a provenance card recording the knowledge assets, reasoning chain, computation steps, and verification results.
\end{itemize}

We evaluate KnowFeat on seven public benchmarks from OpenML and SimECNY, a real-world AML dataset with 200{,}000 digital currency transactions and 93 baseline features. Key findings:

\begin{itemize}
    \item On public datasets, KnowFeat achieves the best AUC on 4 of 7 datasets and the best average rank (1.4) among seven methods, with significantly fewer features than LLM-based alternatives.
    \item On SimECNY, 18 knowledge-guided features improve Recall by +2.1pp and F1 by +1.5pp over the 93-feature baseline, outperforming all baselines.
    \item Ablation confirms that domain knowledge is essential: the schema-only variant generates zero accepted features on one dataset.
\end{itemize}
